\chapter{Glossary}
\label{chap:glossary}
\label{sec:glossary}
This glossary defines the canonical terms used in this manual. Terms are normative unless marked as implementation-defined.

\begin{longtable}{@{}L{0.22\linewidth}L{0.68\linewidth}@{}}
\toprule
Term & Definition \\
\midrule
\texttt{T*} & CPU-visible tile instruction. A tile instruction with a catalog opcode in the instruction stream. Executed by the Davinci front end as a PTO bundled descriptor instruction. Its implementation is a UOP Engine-generated microcode body. Also called a \textit{tile instruction}. Examples: \texttt{TADD}, \texttt{TROWSUM}. \\
\texttt{u\_vec.*} & UOP Engine-generated vector micro instruction. Not directly visible to the CPU instruction stream. Full-tile Vec bodies issue through Vector RS and drive VEC-4K-v2 beat words. Examples: \texttt{u\_vec.vadd}, \texttt{u\_vec.vlds}, \texttt{u\_vec.vsts}. \\
\texttt{V*} & Davinci VTG vector micro-instruction profile. These are implementation micro-ops for 256 B or 512 B VTG execution and are scheduled by GVIQ. They are not PTO \texttt{B.*} descriptor encodings. \\
\textbf{Tile} & A two-dimensional array of elements stored in on-chip tile storage. The fundamental data object of the PTO ISA. Described by tile metadata (shape, valid shape, layout, element format, location intent). \\
\textbf{Tile instruction} & See \texttt{T*}. \\
\textbf{Scalar register} & One of \texttt{X0}--\texttt{X31}. Scalar register fields are 5-bit values; \texttt{X0} is hardwired to zero. \\
\textbf{Tile operand selector} & Encoded tile operand field in a \texttt{B.IOT} descriptor. The opcode profile defines whether the selector names a destination, source, mask, or descriptor-local operand binding. \\
\textbf{Tile metadata table} & Implementation table holding per-tile metadata such as shape, valid rows, valid columns, layout, format, and readiness. In Davinci-v2 this role is carried by the Tile Metadata RAT; PTO rows map to \texttt{shape.y}, columns to \texttt{shape.x}, and element format to \texttt{format}. \\
\textbf{Ready Table} & Implementation bit vector tracking whether selected tile payload state has a valid result. Used by the issue queue to determine when a tile instruction is ready to issue. \\
\textbf{Microcode body} & The sequence of 64-bit micro instruction sites that implement a tile instruction. It is generated by the UOP Engine from the decoded descriptor bundle and implementation templates; it is not fetched as CPU instruction stream bytes. \\
\textbf{WO/WM pair} & One 64-bit static micro instruction site. Bits \texttt{63:32} hold \texttt{WO}, the operation half; bits \texttt{31:0} hold \texttt{WM}, the metadata half. This 8-byte word is the fundamental unit of microcode storage in this manual. \\
\textbf{Beat word} & The expanded vector execution word produced by the microassembler from a 64-bit WO/WM micro instruction and decoded metadata. One static site may expand to multiple beat words, for example for striping or masking. \\
\textbf{Static site} & One 64-bit WO/WM micro instruction in a microcode body. Each static site appears once in the UOP Engine template; loops cause the site to execute multiple times at runtime. The WO/WM sites column in the opcode map counts static sites. \\
\textbf{Vector RS} & Vector Reservation Station. Davinci-v2 issue queue for full-tile VEC-4K-v2 arithmetic, logical, comparison, reduction, select, conversion, and special-function tile bodies. \\
\textbf{GVIQ} & Grouped Vector Issue Queue. Davinci-v2 issue queue for VTG micro-instructions. It uses \texttt{block\_id} and \texttt{pc\_index} to read the micro-instruction buffer and arbitrates with Vector RS for the VEC datapath. \\
\textbf{MTE RS} & Memory Tile Engine reservation station. Davinci-v2 issue queue for tile load, store, gather, scatter, get, put, move, and transpose profiles. \\
\textbf{BROB} & Block Reorder Buffer. Davinci-v2 block-granularity retirement structure used for precise exception reporting, BID-ordered flushing, and side-effect commit. \\
\textbf{BID} & Block ID assigned by Davinci-v2 at block start. It orders BROB entries and lets the backend invalidate younger IQ, GVIQ, SSB, STQ, and rename state on exceptions or flushes. \\
\textbf{Retirement state} & Backend state linking decoded PTO work to block completion and side-effect commit. In Davinci-v2 this is BROB/BID state plus branch tags, SSB/STQ state, and completion signals from Vector RS, MTE RS, Cube RS, or GVIQ. \\
\textbf{Tile buffer} & On-chip storage for tile data. In Davinci-v2 this is TRegFile-4K physical tile storage addressed by implementation tags such as \texttt{PT0}--\texttt{PT255}; those tags are not PTO architectural operand names. \\
\textbf{TRegFile-4K} & Davinci-v2 physical tile register file. Vector, cube, MTE, and VTG execution read and write this storage through physical tile tags after Tile RAT rename. \\
\textbf{VEC datapath} & Davinci vector execution datapath. Consumes SOP or micro-instruction-buffer beat words, reads SA/SB/SC staging data sourced from TRegFile-4K, performs the operation, and writes results through TRegFile-4K writeback paths. \\
\textbf{Staging register} & Davinci VEC-4K-v2 staging state such as \texttt{SA}, \texttt{SB}, \texttt{SC}, and \texttt{SOP}. Staging registers hold fetched tile payloads, predicates, or beat words during execution; they are not PTO architectural registers. \\
\textbf{Mask / predicate} & A per-lane control bit determining whether a lane produces a result. Lanes with mask=0 are inactive and do not participate in arithmetic. Micro instructions operate under a mask emitted by the tile metadata path. \\
\textbf{Valid region} & The subset of a tile's physical shape that is semantically active. Defined by valid rows and valid columns. Operations are defined only over the valid region. Values outside the valid region are implementation-defined. \\
\textbf{Layout} & The memory order interpretation of tile elements. Row-major: consecutive columns share the same row. Column-major: consecutive rows share the same column. Selected by the tile opcode profile or a \texttt{B.ATTR} refinement. \\
\textbf{Shape selector} & Encoded dimension binding carried by \texttt{B.DIM} or by a profile-defined metadata descriptor. Selects the physical tile dimensions and valid region bounds from metadata state. \\
\textbf{Opcode field} & The operation selector in \texttt{B.OP}, exposed in this manual as \texttt{B.OP.Opcode}. \\
\textbf{Microassembler} & The hardware unit that combines a 64-bit WO/WM micro instruction with decoded metadata to produce vector beat words. Also performs variant selection and port recipe application. \\
\textbf{D1 / D2 / D3} & Decode and rename stages in the Davinci front end. D1 performs domain and opcode decode; D2 performs scalar and tile rename; D3 prepares dispatch, readiness, metadata, and issue-queue routing. \\
\textbf{UOP Engine entry table} & Implementation table selected by the decoded tile opcode profile. Maps each entry ID to a body base or generation recipe, body length, variant table pointer, and port recipe pointer. \\
\textbf{B.OP} & Required first descriptor of every bundled tile instruction. It carries the tile opcode, data type, and profile selector. \\
\textbf{Descriptor bundle} & A CPU-visible tile instruction made of \texttt{B.OP} plus the following \texttt{B.*} descriptors required by the selected tile opcode profile. The bundle decodes as one PTO architectural instruction; a Davinci profile may retire it as part of a BROB block. \\
\bottomrule


\textbf{GEMM} & General matrix multiply. The operation $C += A \times B$ where $A$, $B$, and $C$ are matrices. In PTO, GEMM is typically implemented as a nested loop over tiles: \texttt{TMUL} computes one rank-1 block product; \texttt{TADD} or \texttt{TADDS} accumulates it into $C$. \\
\textbf{SPMD} & Single Program Multiple Data. A parallel execution model where all execution agents run the same program but on different data. PTO tile programs are typically written in SPMD style. \\
\textbf{MPMD} & Multiple Program Multiple Data. A parallel execution model where different execution agents run different programs. PTO supports MPMD dispatch alongside SPMD. \\
\textbf{Happens-before} & The ordering guarantee between two operations. If operation A happens-before operation B, the effects of A are visible to B. PTO establishes happens-before ordering through events, \texttt{TSYNC}, or explicit memory dependencies. \\
\textbf{ReLU} & Rectified linear unit. The activation function $\max(0, a)$. In PTO, \texttt{TRELU} implements this as a per-element tile instruction. \\
\textbf{Leaky ReLU} & The activation function $\alpha a$ for $a < 0$, $a$ for $a \ge 0$, where $\alpha$ is a small positive constant (typically $0.01$). In PTO, \texttt{TLRELU} implements this. The slope $\alpha$ is implementation-defined, not an explicit operand. \\
\textbf{Argmax / Argmin} & The index of the maximum or minimum element along an axis. \texttt{TROWARGMAX} returns the column index of the maximum element in each row; \texttt{TROWARGMIN} does the same for the minimum. These are the canonical tile instructions for softmax normalisation and attention kernels. \\
\textbf{Softmax} & The normalisation operation $\mathrm{softmax}(x)_i = \exp(x_i) / \sum_j \exp(x_j)$. Typically implemented in PTO as: subtract the row maximum (\texttt{TROWSUB}), exponentiate (\texttt{TEXP}), sum the row (\texttt{TROWSUM}), then divide each element (\texttt{TDIV}). \\
\textbf{Mixed-precision} & A pipeline that uses different numeric formats at different stages (e.g. f32 accumulation followed by f16 storage). \texttt{TCVT} is the primitive for converting between formats within a mixed-precision pipeline. \\
\textbf{Rank-1 update} & A matrix update of the form $C += u v^\top$ where $u$ is a column vector and $v^\top$ is a row vector. Each outer product produces a full matrix; the \texttt{TMUL} tile instruction computes one block of this product. GEMM decomposes into a sequence of rank-1 updates tiled across the $A$ and $B$ matrices. \\
\textbf{Accumulator tile} & A destination tile that accumulates results from multiple operations without intermediate stores. In PTO GEMM, the accumulator tile starts at zero and is updated in-place by \texttt{TADD} after each rank-1 block product. \\
\textbf{Ragged tile} & A tile whose valid region does not fill its physical shape. This typically occurs at batch edges when the batch size is not evenly divisible by the tile height. Partial operations (\texttt{TPARTADD}, \texttt{TPARTMUL}, etc.) and \texttt{TFILLPAD} handle ragged tiles without requiring the full tile to be padded first. \\
\textbf{LR/SC} & Load-reserved / store-conditional. A pair of instructions that implement atomic compare-and-swap in PTO's scalar ISA. \texttt{lr.d} loads a value and sets a reservation; \texttt{sc.d} stores only if the reservation is still valid. A failed \texttt{sc} returns 1 and indicates that another hart modified the memory location. \\
\textbf{Acquire / Release} & Memory ordering semantics for atomic operations. Acquire (\texttt{.aq}) prevents subsequent loads from moving before the atomic. Release (\texttt{.rl}) prevents prior stores from moving after the atomic. \texttt{.aqrl} provides both. These semantics are needed for correct synchronisation in concurrent programs. \\
\textbf{FMA} & Fused multiply-add. Computes $a \times b + c$ with a single rounding step. In PTO's scalar FP ISA, \texttt{fmadd}, \texttt{fmsub}, \texttt{fnmadd}, and \texttt{fnmsub} provide FMA. FMA gives higher accuracy and better throughput than separate multiply and add. \\
\textbf{FPCR} & Floating-Point Control Register. A CSR that holds the current rounding mode and exception flags. Accessed in PTO through \texttt{fsr}, \texttt{fset}, and \texttt{fget} instructions. The rounding mode in FPCR affects all FP operations that round, unless the instruction overrides it with a variant field. \\
\textbf{IEEE 754} & The international standard for floating-point arithmetic. PTO's scalar and tile FP instructions follow IEEE 754: infinities, NaN values, signed zero, and subnormal numbers all behave as defined by the standard. NaN propagates with its sign bit preserved; division by zero produces Inf with the dividend's sign. \\
\textbf{NaN propagation} & The rule that any arithmetic operation on a NaN operand produces a NaN result. PTO tile instructions follow this: if any lane holds a NaN, the corresponding result lane is a NaN. The sign bit of the input NaN is preserved through the operation. \\
\textbf{Staging register} & A temporary register used within a microcode body to hold an intermediate value between micro-instruction sites. In the \texttt{u\_vec.*} ISA, a \texttt{vec\_stg} source type references a staging register written by an earlier site in the same microcode body. \\
\textbf{Beat word} & The expanded vector execution word produced by the microassembler. One 64-bit WO/WM microcode site can expand to multiple beat words at runtime, for example when the microassembler unrolls a row loop over multiple VEC datapath cycles. The beat word drives the actual hardware ALUs and memory pipelines. \\
\textbf{Static site} & One 64-bit WO/WM micro-instruction entry in a microcode body. The term \emph{static} distinguishes it from the \emph{dynamic} executions that occur when the microassembler unrolls a loop or stripes over lanes. The microcode-site count in the opcode map counts static sites. \\
\textbf{UOP Engine} & A table indexed by the tile opcode/profile that maps each tile instruction to its microcode body. Each directory entry points to a body base address, body length, variant table, and port recipe. The microassembler uses the directory to locate and expand the body. \\
\textbf{Microassembler} & The hardware unit that converts one 64-bit WO/WM static site plus decoded tile metadata into one or more VEC beat words. The microassembler handles loop unrolling (producing multiple dynamic executions from one static site), lane striping, mask application, and variant selection. \\
\textbf{Micro-instruction buffer} & A hardware cache that stores hot microcode bodies near the execution pipeline. Some implementations store the most frequently used microcode bodies here to avoid the UOP Engine microcode generation on common tile instruction sequences. The buffer is an implementation optimisation, not an architectural requirement. \\
\textbf{Vector RS} & Vector Reservation Station. The issue queue that holds full-tile VEC arithmetic, logical, comparison, reduction, select, conversion, and special-function operations until their operands are ready. In Davinci-v2, Vector RS drives VEC-4K-v2 execution. \\
\textbf{MTE RS} & Memory Tile Engine Reservation Station. The issue queue that holds tile load, store, gather, scatter, get, put, move, and transpose operations. MTE RS manages the memory pipeline and interfaces to the tile storage subsystem. \\
\textbf{GVIQ} & Grouped Vector Issue Queue. In Davinci-v2, the issue queue for VTG vector micro-instructions. GVIQ holds 256 B or 512 B VTG micro-ops, uses \texttt{block\_id} and \texttt{pc\_index} to read the micro-instruction buffer, and arbitrates with Vector RS for the VEC datapath. \\
\textbf{BROB} & Block Reorder Buffer. The Davinci-v2 block-granularity retirement structure. It provides precise exception reporting, BID-ordered flushing, and side-effect commit for block-structured tile programs. BROB entries are ordered by Block ID. \\
\textbf{BID} & Block ID. A unique identifier assigned by Davinci-v2 at block start. BID orders BROB entries and enables the backend to invalidate younger IQ, GVIQ, SSB, STQ, and rename state on exceptions or flushes. \\
\textbf{Tile RAT} & Tile Register Alias Table. The rename table that maps architectural tile operand selectors to physical tile tags. In Davinci-v2, tile operands go through rename; the Tile RAT checkpoints state for recovery on exceptions and flushes. \\
\textbf{TRegFile-4K} & Davinci-v2 physical tile register file. All vector, cube, MTE, and VTG execution reads and writes this storage through physical tile tags after rename. The 4K capacity refers to the total element storage, not the number of tile slots. Physical tags (\texttt{PT0}--\texttt{PT255}) are implementation names, not PTO architectural identifiers. \\
\textbf{Tile Metadata RAT} & The rename table that supplies per-tile metadata: shape (\texttt{shape.y}, \texttt{shape.x}), valid rows, valid columns, layout, element format, and readiness. PTO rows map to \texttt{shape.y}, columns to \texttt{shape.x}, and element format to the Davinci \texttt{format} field. \\
\textbf{D1 / D2 / D3} & Decode and rename stages in the Davinci front end. D1 performs domain and opcode decode. D2 performs scalar and tile rename. D3 prepares dispatch, readiness, metadata, and issue-queue routing. \\
\textbf{PC-relative} & An addressing mode where the effective address is computed as PC + immediate offset. PTO uses PC-relative offsets for branches (\texttt{b.*}, 16-bit signed offset) and for PCR-relative loads (\texttt{l*.pcr}). This makes code position-independent without requiring a global pointer register. \\
\textbf{PC-relative load} & A load instruction that computes its address as PC + immediate offset. In PTO: \texttt{lb.pcr}, \texttt{lh.pcr}, \texttt{lw.pcr}, \texttt{ld.pcr}. Used for accessing small, position-independent data objects embedded near code, such as jump tables and small constants. Unlike base+offset loads, PC-relative loads do not require a base register to be set first. \\
\textbf{Position-independent code} & Code that executes correctly regardless of where it is loaded in memory. PTO tile programs are inherently position-independent: tile instructions are PC-relative or use tiles as operands rather than absolute memory addresses. PCR-relative loads (\texttt{l*.pcr}) further support PIC scalar code. \\
\textbf{SSA form} & Static Single Assignment form. An IR representation where each variable is assigned exactly once. In the compiler backend, SSA separates integer and FP register uses, so type confusion between integer and FP values held in the same \texttt{X0}--\texttt{X31} register file is a frontend concern, not a backend concern. \\
\textbf{Register allocation} & The compiler pass that assigns infinite SSA values to the finite set of physical registers \texttt{X0}--\texttt{X31}. PTO has one flat scalar register file shared between integers and FP values; the allocator must track which register holds which type at any point. \\
\textbf{VTG} & Vector Thread Generator. The Davinci-v2 component that emits VTG micro-instructions from cached microcode bodies or the micro-instruction buffer. VTG micro-instructions (\texttt{V*}) are distinct from \texttt{u\_vec.*} micro-instructions; they share the same VEC datapath but use a different issue surface (GVIQ). \\
\end{longtable}
\index{glossary}
\index{GEMM}\index{SPMD}\index{MPMD}\index{ReLU}\index{FMA}\index{FPCR}\index{IEEE 754}\index{NaN propagation}\index{Acquisition}\index{Release ordering}\index{LR/SC}\index{GVIQ}\index{MTE RS}\index{TRegFile-4K}\index{Tile RAT}\index{BROB}\index{BID}\index{D1}\index{D2}\index{D3}\index{PC-relative}\index{Position-independent code}\index{SSA form}\index{Register allocation}\index{Beat word}\index{Static site}\index{Microassembler}\index{Micro-instruction buffer}\index{Rank-1 update}\index{Accumulator tile}\index{Ragged tile}\index{Staging register}\index{Softmax}\index{Mixed-precision}\index{Argmax}\index{Argmin}\index{Leaky ReLU}
\index{glossary}
